\pagestyle{plain} % No headers, just page numbers
% \pagenumbering{arabic} % Arabic numerals
% \setcounter{page}{1}


\chapter{\uppercase {Human-Inspired Jailbreak Defenses }}


Understanding malicious intent extends beyond detecting passive deception to actively defending against adversarial attacks. Jailbreak attacks on large language models represent a form of intentional manipulation where adversarially crafted prompts attempt to circumvent safety guardrails and elicit harmful or prohibited content. These attacks succeed by exploiting various vulnerabilities, including indirect phrasing, role-playing scenarios, multi-turn persuasion, and prompt injection techniques. The problem is particularly acute because attackers continuously adapt their strategies, creating an arms race between attack sophistication and defense mechanisms. Current defense approaches face significant limitations, including high computational costs, increased latency, reliance on additional LLM calls for input filtering or output verification, and brittleness against novel attack strategies. Moreover, many defenses prove inadequate against sophisticated multi-turn jailbreaks where human adversaries gradually manipulate the conversation state to bypass restrictions. Recent research ~\cite{mhj} demonstrates that multi-turn human jailbreaks uncover significant vulnerabilities, exceeding 70\% attack success rate (ASR) against defenses that report single-digit ASRs with automated single-turn attacks.

Just as deception detection requires recognizing when speakers deliberately mislead, jailbreak defense demands identifying when users intentionally attempt to manipulate the model into violating its safety policies. The human strategies we studied for detecting scams and phishing attacks provide valuable insights into how people recognize and resist social engineering, which shares fundamental characteristics with jailbreak attempts. Both involve deceptive framing, gradual persuasion, authority exploitation, and emotional manipulation. By treating jailbreak defense as a form of malicious intent detection and drawing inspiration from human resistance strategies, we propose a novel approach that enhances the model's intrinsic ability to recognize and refuse adversarial requests through explicit reasoning that mirrors how humans identify and respond to manipulation attempts.

\section{Preliminaries}

\subsection{Jailbreak Attack Techniques}

Jailbreak attacks employ diverse strategies to manipulate LLMs into generating prohibited content. \textbf{Gradient-based attacks} such as Greedy Coordinate Gradient (GCG) \cite{zou2023universal} optimize adversarial suffixes by computing gradients with respect to the model's token embeddings, generating universal strings that can be appended to harmful queries to induce compliance. GCG demonstrates remarkable transferability across different models and has inspired numerous variants that improve efficiency and effectiveness. However, these automated attacks often produce nonsensical text patterns that can be detected through perplexity-based filtering.

\textbf{Prompt-based attacks} leverage carefully crafted instructions without gradient optimization. Prompt Automatic Iterative Refinement (PAIR) \cite{chao2023jailbreaking} uses an attacker LLM to iteratively refine jailbreak prompts based on the target model's responses, achieving high success rates through automated red-teaming. These methods exploit the model's instruction-following capabilities by embedding harmful requests within complex scenarios, role-playing contexts, or multi-step reasoning chains. The Jailbreak Benchmark (JBB) \cite{chao2024jailbreakbench} consolidates 100 challenging behaviors spanning disinformation, illegal activities, privacy violations, and other harmful categories to systematically evaluate attack success rates.

\textbf{Multi-turn attacks} represent the most challenging threat vector. Unlike single-shot attempts, multi-turn human jailbreaks (MHJ) \cite{mhj} involve sustained conversation where attackers gradually build context, establish rapport, or incrementally escalate requests across multiple exchanges. Research shows that even models resistant to single-turn attacks succumb to multi-turn persuasion at alarming rates, with success exceeding 70 percent in some cases \cite{m2s}. These attacks exploit the model's tendency to maintain conversational coherence and its difficulty in tracking malicious intent across extended interactions.

\subsection{Existing Defense Mechanisms}

Defense strategies broadly fall into three categories. \textbf{Input filtering approaches} attempt to detect and block malicious prompts before processing. LlamaGuard \cite{llamaguard} employs a specialized LLM fine-tuned on safety taxonomies to classify input prompts as safe or unsafe across multiple risk categories. While effective against known attack patterns, these methods incur significant latency by requiring an additional LLM inference before each response. They also struggle with novel attack formulations and can be circumvented through encoding schemes or indirect phrasing.

\textbf{Output filtering mechanisms} inspect generated responses for harmful content. SafeDecoding \cite{xu2024safedecoding} modifies the decoding process to reduce the probability of unsafe tokens by incorporating safety classifiers into beam search. Erase-and-Check \cite{kumar2023certified} generates multiple candidate responses, identifies potentially harmful segments, removes them, and verifies the remaining content's coherence and safety. These approaches face challenges in distinguishing legitimately sensitive information from genuinely harmful content, leading to either over-censorship or exploitation through carefully worded requests.

\textbf{Adversarial robustness techniques} aim to make models inherently resistant to attacks. SmoothLLM \cite{robey2023smoothllm} adds random perturbations to input prompts and aggregates predictions across multiple perturbed versions, effectively neutralizing gradient-based attacks like GCG. However, this defense requires multiple forward passes, substantially increasing computational costs. PerplexityFilter \cite{alon2023detecting} detects adversarial suffixes by measuring input perplexity. While effective against optimization-based attacks, this method fails against human-crafted jailbreaks that maintain natural language characteristics.

Recent work explores self-defensive capabilities within LLMs themselves. Self-Reminder \cite{xie2023defending} prepends system messages reminding the model of safety policies before processing each request. Although simple and efficient, this approach proves vulnerable to advanced attacks that explicitly acknowledge then override these reminders. SelfDefend \cite{wang2024selfdefend} enables models to first assess whether a prompt violates safety policies before generating a response, implementing a form of deliberative refusal. This method demonstrates improved robustness while maintaining low latency but relies heavily on the quality of safety reasoning learned during training.

\subsection{Limitations of Current Approaches}

The practical deployment of jailbreak defenses must satisfy four critical criteria as identified by~\cite{wang2024selfdefend}. First, defenses must provide \textbf{strong protection} against diverse attack strategies including both automated and human-crafted attempts across single-turn and multi-turn scenarios. Second, they must maintain \textbf{minimal latency} to ensure acceptable user experience, ruling out approaches requiring multiple LLM calls or extensive post-processing. Third, defenses should demonstrate \textbf{broad generalization} to novel attack formulations without requiring frequent updates or retraining. Fourth, they must preserve \textbf{high utility} by avoiding false positives that incorrectly refuse legitimate queries or degrade response quality.

Evaluating existing methods against these criteria reveals significant gaps. Input and output filtering approaches using secondary LLMs fundamentally violate the latency requirement, doubling inference time and computational costs. LlamaGuard and similar classifiers add significant time overhead per request depending on model size, making them impractical for latency-sensitive applications. These methods also struggle with generalization, as new attack patterns regularly emerge that bypass static classification rules. The adversarial brittleness problem is particularly severe, attackers can iteratively probe filter responses to identify boundary cases where harmful content slips through.

Methods claiming human-inspired reasoning often fall short of genuinely mimicking human defensive strategies. SafeBehavior \cite{guan2024safebehavior} proposes multi-stage reasoning that superficially resembles human deliberation but implements this through rigid rule-based chains that lack the adaptive flexibility characterizing human judgment. The approach generates extensive reasoning traces analyzing potential harms across predefined categories, but this process is both computationally expensive and vulnerable to attacks that acknowledge each stage's concerns before providing counter-arguments. Similarly, recent work on constitutional AI \cite{bai2022constitutional} incorporates ethical principles into model responses through lengthy chain-of-thought generation, but the resulting explanations often follow templated patterns rather than exhibiting the nuanced context-dependent reasoning humans employ when detecting manipulation.

The fundamental limitation is that current defenses lack the adaptive pattern recognition humans develop through experience with deception. When humans encounter phishing emails, scam calls, or social engineering attempts, they draw upon learned heuristics, pattern matching from previous exposures, and contextual reasoning about the requester's true intentions. They recognize red flags such as unusual urgency, requests for sensitive information, authority impersonation, and inconsistencies in the narrative. Crucially, humans maintain suspicion across conversation turns, remembering earlier red flags even when subsequent messages seem innocuous. Existing LLM defenses typically evaluate each prompt in isolation or through shallow contextual analysis, failing to accumulate suspicion scores or recognize manipulation tactics that unfold gradually.

Multi-turn attacks expose these weaknesses most clearly. In the MHJ dataset \cite{mhj}, attackers often begin with seemingly benign questions to establish rapport, then introduce policy-violating elements incrementally, and finally escalate to explicit harmful requests once the model's guard is lowered. Current defenses struggle because they either assess each turn independently, losing critical context about the attacker's evolving strategy, or attempt to maintain comprehensive conversation history which becomes computationally prohibitive and still fails to properly weight early warning signals. Human defenders naturally track conversation trajectories and raise their suspicion threshold when they detect manipulative patterns, a capability largely absent in automated defenses.

\section{Proposed Approach: Human-Inspired Defense Training}

We propose enhancing LLMs' intrinsic jailbreak resistance by training them to recognize and respond to malicious requests using strategies inspired by how humans detect and resist real-world deception. Our approach consists of three interconnected components addressing data collection, model training, and evaluation.

\subsection{Human Deception Defense Data Collection}

We plan to construct a diverse dataset capturing human strategies for recognizing and resisting manipulation across multiple domains. First, we will systematically collect examples from online forums where users share experiences with scam attempts. Communities such as r/Scams on Reddit contain thousands of detailed accounts where people describe attempted frauds, explain what made them suspicious, and discuss how they verified legitimacy or refused compliance. We will extract structured examples capturing the deceptive request, the red flags identified, and the defensive response strategy. For instance, a typical example might document a phishing email requesting password verification, note the suspicious sender domain and unusual urgency language as red flags, and describe the defensive action of directly contacting the organization through official channels rather than clicking provided links.

Second, we will systematically analyze existing jailbreak attack datasets to extract implicit human defensive reasoning. For each attack prompt in datasets like JBB, MaliciousInstruct, and MHJ, we will use a combination of expert annotation and LLM-assisted analysis to generate explanations of why the request violates policies and how a human would recognize the manipulation attempt. This creates training examples where deceptive requests are paired with explicit, reasonable reasoning chains identifying the malicious intent, the policy violation category, and the appropriate refusal strategy.

Critically, we will ensure the dataset includes extensive multi-turn scenarios where manipulation unfolds gradually. For each multi-turn attack from the MHJ dataset, we will annotate each conversation turn with cumulative suspicion scores indicating how human defenders would adjust their vigilance as the conversation progresses. This addresses the key weakness in existing defenses by providing training signal for maintaining and accumulating suspicion across turns rather than evaluating each message in isolation.

\subsection{Training Methodology}

Our training approach combines supervised fine-tuning and reinforcement learning to instill human-like defensive reasoning. In the supervised fine-tuning phase, we will train the model on our collected dataset using a specific format that encourages explicit reasoning before responding. Each training example will follow the pattern: given an input request (potentially within multi-turn context), the model first generates a brief internal reasoning trace assessing potential red flags, then produces the appropriate response (either helpful compliance for benign requests or polite refusal with explanation for malicious ones).

The reasoning trace will be structured to mirror human defensive cognition. For example, when encountering a request to provide instructions for illegal activities, the reasoning might identify policy violation categories, note absence of legitimate use cases, recognize direct vs. indirect phrasing, and decide on refusal strategy. Importantly, for multi-turn scenarios, the reasoning explicitly references earlier conversation turns, accumulates suspicion signals, and adjusts the refusal threshold accordingly. This contrasts sharply with existing methods that generate reasoning only when refusing, we will train the model to briefly reason about safety for all requests, building a habit of vigilance.

Following supervised fine-tuning, we will employ Group Relative Policy Optimization (GRPO) \cite{shao2024deepseekmath} for reinforcement learning with carefully designed verifiable rewards. The reward function assesses whether the model successfully resisted jailbreak attempts while maintaining helpfulness on benign queries. For malicious requests, positive reward requires both correct refusal and presence of valid reasoning identifying the manipulation tactics. For multi-turn attacks, we implement a cumulative reward structure penalizing compliance at any point in the conversation and rewarding consistent refusal that references earlier red flags. Conversely, the model receives positive reward for helpfully responding to legitimate queries and negative reward for false refusals that cite non-existent safety concerns. The verifiable nature of jailbreak defense outcomes (did the model provide harmful content or not) allows us to construct deterministic reward signals rather than relying on expensive human preference judgments. We will iteratively sample multiple responses for each prompt, compute rewards, and update the model to increase the probability of high-reward defensive behaviors.

This human-inspired approach addresses the core limitations of existing defenses by enhancing the model's intrinsic understanding of malicious intent rather than relying on brittle external filters or computationally expensive multi-step verification. By learning from how humans recognize and resist deception through experience and adaptive pattern matching, we aim to develop defenses that are simultaneously more robust, more efficient, and more naturally integrated with the model's reasoning capabilities. This work represents the natural culmination of our thesis's focus on intent understanding, extending from benign communication goals through deception detection to adversarial manipulation resistance as different manifestations of the same fundamental challenge.